% ============================================================
% 问题3 小节（国赛风格）—— 插入 001sec.txt 问题2 小节之后。
% 数据口径：3/A3-codex/A3-codex（tight800 最终交付，与问题2 A2-gpt 同源）。
% 主结果：末值维持口径（与4.2一致）连续临界 57.169466 h、
%         首次严格达标整分钟 57.1833 h（205860 s）。
% 图片：fig4-3-1.png（含水率时空分布）、fig4-3-2.png（径向剖面）、
%       fig4-3-3.png（干燥曲线）、fig4-3-4.png（外推敏感性），
%       取自 A3-codex/figures/01~04，复制到论文图片目录后按此名引用。
% ============================================================
\subsection{问题3：满足烘干要求的最短烘干时间}

\subsubsection{问题分析}

按烘干要求，药材\textbf{各处}水分浓度应低于 $0.15\ \mathrm{kg/kg}$。数值上两问的“最大
含水率”始终位于中心（表面先干、中心最后干，后文给出数值印证），故判据等价于中心判据，
烘干时间可写成事件时刻
\begin{equation}
  t^*=\min\big\{\,t:\ \max_r C(r,t)<0.15\ \mathrm{kg/kg}\,\big\}.
  \label{eq:event}
\end{equation}
本问是问题 2 模型的全流程版本：控制方程、本构（附录 3）、边界条件（附件 1）与数值框架
均与 4.2 节相同，唯一区别是考察窗由固定 3 h 变为由判据事件终止的长时间积分。

先做量级判断。水分柱对称基模时间常数 $\tau_1=R^2/(\lambda_1^2D_0)\approx3.4\ \mathrm{h}$；
若扩散系数恒为初值，基模衰减 $2.55\to0.15$ 仅需
$\tau_1\ln(2.55/0.15)\approx9.6\ \mathrm{h}$。而由 $D\propto e^{-0.45/C}$，$C$ 从 2.55
降至 0.15 时浓度因子衰减 16.8 倍（温度项在恒温段仅给出约 2.4 倍加速，净衰减约 7 倍），
$\tau_1$ 随之增至约 24 h——后期“越烘越慢”的自锁效应把过程拉长约 6 倍，题面“2–3 天”的
干燥周期正是其宏观表现。

\subsubsection{模型建立}

控制方程与定解条件同式~\eqref{eq:heat2}--\eqref{eq:propsD2}、\eqref{eq:bcR2}：
物性按附录 3 取 $\rho=650+128C$、$c_p=1450+2736C/(C+1)$、$k=0.21+0.38C/(C+1)$、
$D=2.4\times10^{-3}\exp(-0.45/C)\exp(-3850/T)$；初值 $T=28\,^{\circ}\mathrm{C}$、
$C=2.55\ \mathrm{kg/kg}$；烘房条件由附件 1 时变序列给定（相邻观测点线性插值），
$14400\,\mathrm{s}$ 后\textbf{按末值维持}（$T_a=50.1650\,^{\circ}\mathrm{C}$、
$C_a=0.04986\ \mathrm{kg/kg}$）——这是全流程唯一需要外推的假设，其影响在模型讨论中量化。
半径取 $R=2\ \mathrm{cm}$ 常数，收缩由问题 4 处理。

\subsubsection{模型求解}

沿用 4.2 节的单元中心有限体积框架并作两点强化：表面加密网格
$\xi_{j+1/2}=1-(1-j/N)^2$（$N=800$），界面系数沿相邻状态线性路径作三点 Gauss 积分，真实
表面由耦合不动点迭代解出（残差阈值 $2\times10^{-13}$）；时间积分取 BDF
（`scipy.solve_ivp`，$\mathrm{rtol}=10^{-11}$、$\mathrm{atol}_T=10^{-12}$、
$\mathrm{atol}_C=10^{-13}$，块三对角稀疏 Jacobian），前 4 h 每 60 s 数据节点重启
（$\max\Delta t=10\ \mathrm{s}$），此后连续积分（$\max\Delta t=300\ \mathrm{s}$）。
烘干时间以终止型事件函数 $g(t)=\max_r C(r,t)-0.15$ 直接根搜索（式~\eqref{eq:event}），
由 BDF 积分器收缩步长做 Brent 定位，不引入后处理插值误差；事件函数取按 $r^2$ 二次重构的
中心值、全部单元值与真实表面值的最大者。先搜索 72 h，未达标按 24 h 延长，30 天仍未穿越
则报错而不伪造时长。

输出按题面要求：result3.xlsx 自 $t=60\,\mathrm{s}$ 起每隔 60 s、到中心距离每隔 0.1 cm
（$0\sim1.9\ \mathrm{cm}$ + 表面列），至首次严格达标整分钟为止，共 3431 行、四位小数。

\subsubsection{结果与分析}

烘干时间的求解结果：连续临界时刻为 \textbf{57.169466 h}（此时最大含水率恰等于 0.15，
严格低于阈值需略晚）；按题目 60 s 采样精度，首次全场严格达标为 \textbf{57.1833 h，
即 57 小时 11 分钟（205860 s）}。论文可写约 57.18 h（2.38 天），与题面“2–3 天”一致。
表~\ref{tab:C5} 给出题面表 5 格式的水分浓度。

\begin{table}[htbp]
  \centering
  \caption{表 5\quad 药材烘干过程的水分浓度（$\mathrm{kg/kg}$）}
  \label{tab:C5}
  \begin{tabular}{cccccc}
    \hline
    时间/h & \multicolumn{5}{c}{到药材中心的距离/cm} \\
    & 0 & 0.5 & 1.0 & 1.5 & 2.0 \\
    \hline
    6  & 1.0156 & 0.9868 & 0.9006 & 0.7546 & 0.5335 \\
    12 & 0.4548 & 0.4418 & 0.4019 & 0.3289 & 0.1639 \\
    18 & 0.2979 & 0.2906 & 0.2679 & 0.2247 & 0.0842 \\
    24 & 0.2371 & 0.2321 & 0.2161 & 0.1851 & 0.0663 \\
    30 & 0.2053 & 0.2013 & 0.1887 & 0.1637 & 0.0597 \\
    36 & 0.1854 & 0.1821 & 0.1714 & 0.1500 & 0.0565 \\
    42 & 0.1717 & 0.1687 & 0.1593 & 0.1404 & 0.0547 \\
    48 & 0.1615 & 0.1588 & 0.1504 & 0.1332 & 0.0536 \\
    54 & 0.1535 & 0.1511 & 0.1433 & 0.1275 & 0.0528 \\
    57.1833（烘干结束） & 0.1500 & 0.1477 & 0.1402 & 0.1249 & 0.0525 \\
    \hline
  \end{tabular}
\end{table}

\begin{figure}[H]
  \centering
  \includegraphics[width=0.9\textwidth]{4-3pic/fig4-3-3.png}
  \caption{表面/中心水分浓度全流程时间历程：虚线为 0.15 判据，烘干在 57.18 h 完成}
  \label{fig:drying3}
\end{figure}

\textbf{(1) 温度场全程准稳态。}预热平衡结束时（2 h）表面温度已达 $49.85\,^{\circ}\mathrm{C}$，
3 h 后表面与烘房仅差 $0.23\,^{\circ}\mathrm{C}$；整个恒温干燥段（2–57.2 h）全场温度维持
$50.1650\,^{\circ}\mathrm{C}$、径向温差可忽略，与 $\tau_T\approx46\ \mathrm{min}$ 的预估
一致——热学参数对烘干时间几乎无影响（后文敏感性）。

\textbf{(2) 水分场三阶段。}表~\ref{tab:C5} 中心列给出阶段划分的直接证据：① \textbf{快速失水}
（0–6 h）：中心降幅 $1.5344$，完成中心总降幅的 $63.9\%$，对应表面耗竭与干燥锋推进；
② \textbf{恒温降速干燥}（6–30 h）：每 6 h 中心降幅由 $0.5608$ 递减至 $0.0318$，剖面整体
下移；③ \textbf{末段自锁}（30–57.2 h）：降幅按固定比率衰减（$0.69\sim0.84$），末 3.18 h
中心仅降 $0.0035$，接近指数尾——这是 $D\propto e^{-0.45/C}$ 自锁效应的数学表现：末段
$\tau_1$ 增至约 24 h，松弛极慢。

\textbf{(3) 表面动力学。}表面水分 6 h 已降至 0.5335、12 h 降至 0.1639，此后 45 h 仅再降
0.11（至 0.0525），缓慢逼近烘房平衡值 $C_a=0.04986$（结束差 0.0026，为维持内部出水的
有限通量所需）；中心–表面差由 6 h 的 0.4821 收窄到结束的 0.0975，中心始终是全场最大含水率
位置——“各处低于 0.15”等价于中心判据得到数值印证。

\begin{figure}[H]
  \centering
  \includegraphics[width=0.62\textwidth]{4-3pic/fig4-3-1.png}
  \caption{水分浓度时空演化（57.2 h × 0–2 cm）：虚线为判据前沿 $C=0.15$，
  约 13 h 自表面出发、57.18 h 到达中心}
  \label{fig:moistmap3}
\end{figure}

\subsubsection{模型检验}

\textbf{(1) 网格与容差收敛。}空间 $N=400/800/1600$、容差收紧与 Radau 时间算法对照的结果
见表~\ref{tab:verify3}：含水率观测收敛阶 2.0000；以 800/1600 原始差、容差差与 BDF/Radau
差之和作保守误差估计，含水率 $1.771\times10^{-6}\ \mathrm{kg/kg}$、临界时刻差 $0.239\ \mathrm{s}$，
远小于四位小数阈值 $5\times10^{-5}$，表~\ref{tab:C5} 与 result3.xlsx 的四位小数输出
可靠（末位受舍入分界影响，不作逐位收敛宣称）。

\begin{table}[htbp]
  \centering
  \caption{问题 3 数值对照（共同 60 s × 0.1 cm 输出网格）}
  \label{tab:verify3}
  \begin{tabular}{lccc}
    \hline
    对照 & 含水率最大差/(kg/kg) & 温度最大差/$^{\circ}\mathrm{C}$ & 临界时刻差/s \\
    \hline
    400$\to$800 单元 & $7.074\times10^{-6}$ & $1.388\times10^{-5}$ & 0.9535 \\
    800$\to$1600 单元 & $1.769\times10^{-6}$ & $3.474\times10^{-6}$ & 0.2385 \\
    容差收紧 & $1.042\times10^{-9}$ & $1.404\times10^{-8}$ & 0.0000 \\
    BDF vs Radau & $1.151\times10^{-9}$ & $1.601\times10^{-8}$ & 0.0000 \\
    \hline
  \end{tabular}
\end{table}

\textbf{(2) 守恒与离散恒等式。}水分积分量 $M=\sum V_iC_i$ 满足 $\dot M=-Q_C(R)$：以三点
Gauss 法积分真实边界通量后相对闭合误差 $2.668\times10^{-11}$；显热诊断量
$E=\sum V_iA(C_i)T_i$ 按链式法则 $\dot E=-Q_T(R)+\sum V_iA'(C_i)T_i\dot C_i$（第二项为变
热容的导数修正）闭合相对误差 $3.011\times10^{-10}$；质量与热方程的离散恒等式残差分别为
$6.6\times10^{-24}$、$1.4\times10^{-17}$。这些验证说明所选有效 PDE 被正确离散，不代表
完整蒸发质量—能量模型守恒（见 4.2 节显热近似讨论）。

\textbf{(3) 与问题 2 的衔接。}与问题 2 已保存的 tight800 轨迹逐段核对，0–10800 s 每 60 s
采样最大差为含水率 $5.775\times10^{-12}$、温度 $6.641\times10^{-10}$，第三问准确承接
第二问；全程单元含水率沿径向不增（最大向外增量 $8.9\times10^{-16}$，舍入误差量级），
中心最湿判据成立。

\subsubsection{模型讨论：环境外推的敏感性}

附件 1 只采到 14400 s，其后“维持末值”是全流程最大不确定度。表~\ref{tab:scen3} 给出
外推口径与恒温段扰动的敏感性（N=400，与同网格基准相比）：末小时时间加权均值使临界时刻
推迟 0.31 h；恒温段温度 $\pm1\,^{\circ}\mathrm{C}$ 引起 $\pm1.8\ \mathrm{h}$
（$\mp3.1\%$），湿度 $\pm10\%$ 仅 $\mp0.1\ \mathrm{h}$；按恒温段设定值
$50.0\,^{\circ}\mathrm{C}$ 维持则为 57.48 h（$+0.5\%$，问题 4 同口径并列给出）。机理是
$D$ 的 Arrhenius 因子 $e^{-3850/T}$：温度每升 1 K，$D$ 增大约 $3.7\%$。即便取最不利的
$\pm5\,^{\circ}\mathrm{C}$ 系统偏移，烘干时间区间 $48.9\sim67.6\ \mathrm{h}$ 仍整体落在
“2–3 天”（48–72 h）内，答案对该假设稳健。

\begin{table}[htbp]
  \centering
  \caption{4 h 后烘房条件维持口径对问题 3 烘干时间的影响}
  \label{tab:scen3}
  \begin{tabular}{lcc}
    \hline
    外推情景 & 临界时长/h & 相对同网格基准变化/h \\
    \hline
    末值维持（主口径） & 57.17 & — \\
    恒温段设定值 $50\,^{\circ}\mathrm{C}$ & 57.47 & $+0.30$ \\
    末小时时间加权均值 & 57.48 & $+0.31$ \\
    恒温段温度 $-1\,^{\circ}\mathrm{C}$ & 59.02 & $+1.85$ \\
    恒温段温度 $+1\,^{\circ}\mathrm{C}$ & 55.40 & $-1.77$ \\
    恒温段环境含水率 $-10\%$ & 57.08 & $-0.09$ \\
    恒温段环境含水率 $+10\%$ & 57.28 & $+0.11$ \\
    \hline
  \end{tabular}
\end{table}

综上，固定几何下药材烘干所需时间为 \textbf{57.1833 h（2.38 天）}；完整结果按模板输出至
result3.xlsx（60 s $\times$ 0.1 cm，3431 行，4 位小数），末行中心显示 0.1500 为四位小数
舍入所致（未舍入值 $0.14999\ldots$，严格达标按未舍入值判定）。与问题 4 的收缩模型交叉
（中心水分曲线约 41.0 h 相交、两问干时同落“2–3 天”）见 4.4 节。
